% CVPR 2026 author-kit document setup; camera-ready mode displays all authors.
\documentclass[10pt,twocolumn,letterpaper]{article}
\usepackage{cvpr}
%% This file contains a number of tweaks that are typically applied to the main document.
%% They are not enabled by default, but can be enabled by uncommenting the relevant lines.


%%
%% Official support for 2-column teaser figure
%%
\usepackage{cuted} %< for strip environment

%%
%% For auto-referencing of figures (see fig/teaser.tex and \label{} therein) 
%%
\usepackage{currfile} %< for \currfilebase macro

%%
%% caption and (global) spacing overrides
%%
\usepackage{caption} %< for "\captionof" commands (teaser)
% \setlength{\abovecaptionskip}{.5em} % space above caption
% \setlength{\belowcaptionskip}{1em} % space below caption

%%
%% Inline annotations; for predefined colors, refer to "dvipsnames" in the xcolor package:
%% https://tinyurl.com/overleaf-colors
%%
\newcommand{\red}[1]{{\color{red}#1}}
\newcommand{\todo}[1]{{\color{red}#1}}
\newcommand{\TODO}[1]{\textbf{\color{red}[TODO: #1]}}
%%
%% Disable for camera ready / submission by uncommenting these lines:
%%
% \renewcommand{\TODO}[1]{}
% \renewcommand{\todo}[1]{#1}

%%
%% Work harder in optimizing text layout. Typically shrinks text by 1/6 of page, enable
%% it at the very end of the writing process, when you are just above the page limit.
%%
% \usepackage{microtype}

%%
%% Fine-tune paragraph spacing.
%%
% \renewcommand{\paragraph}[1]{\vspace{.5em}\noindent\textbf{#1.}}

%%
%% Globally adjusts space between figure and caption.
%%
% \setlength{\abovecaptionskip}{.5em}

%%
%% Allows "the use of \paper to refer to the project name"
%% with automatic management of space at the end of the word.
%%
% \usepackage{xspace}
% \newcommand{\paper}{ProjectName\xspace}

%%
%% Commonly used math definitions.
%%
% \DeclareMathOperator*{\argmin}{arg\,min}
% \DeclareMathOperator*{\argmax}{arg\,max}

%%
%% Tighten underline.
%%
% \usepackage{soul}
% \setuldepth{foobar}

%%
%% Reconfigure paragraph (tweaks spacing)
%%

\definecolor{cvprblue}{rgb}{0.21,0.49,0.74}
\usepackage[pagebackref,breaklinks,colorlinks,allcolors=cvprblue]{hyperref}
\def\paperID{}
\def\confName{CVPR}
\def\confYear{2026}
\input{generated/dmd_diversity/settings.tex}
\title{Teacher-Guided Diversity Matching for One-Step Diffusion}
\author{Le Duc Dong \quad Nguyen Huu Thuat \quad Pham Viet Hoang \quad La Nguyen Gia Hy\\CE7454 Project 2}
\begin{document}
\maketitle
\begin{abstract}
We study direct one-step Stable Diffusion distillation under a fixed low-rank
student budget. Starting from the same DMD-inspired checkpoint, we compare
latent and perceptual regression, with and without a teacher-calibrated
same-prompt diversity objective. Four controlled variants share training
captions, teacher pairs, generator updates, and evaluation noise. This draft
imports measured results; one training seed limits the strength of conclusions.
\end{abstract}
\section{Motivation and related work}
Distribution Matching Distillation (DMD) uses real and fake score estimates
to train a direct one-step generator, with perceptual regression against
teacher-generated samples~\cite{yin2024dmd}. Our existing implementation uses
latent MSE and rank-matched LoRA adapters. The question here is whether
perceptual regression and explicit same-prompt variation matching improve this
restricted setting. Cross-sample relational distillation is established in
RDD~\cite{feng2024rdd}; CoDMD studies relational regularization for few-step
video~\cite{zhang2026codmd}. We evaluate a specific image-level extension rather
than claim a new divergence or a reproduction of their training recipes.
\section{Method}
A rank-\DMDStudentRank{} student maps initial Gaussian noise and a text caption
to a clean latent in one denoiser evaluation. The frozen VAE supplies RGB images
for training losses. Teacher and auxiliary score estimates supply the normalized
DMD gradient; the auxiliary scorer is fitted to student-generated noisy latents.
Our objective is
\begin{equation}
L=L_{\mathrm{DMD}}+0.25L_{\mathrm{reg}}+\lambda_d L_{\mathrm{div}}.
\end{equation}
Regression is either latent MSE or LPIPS-VGG~\cite{zhang2018lpips} after decoding
and clamping RGB to $[-1,1]$. Each cached regression group contains two teacher
outputs for the same prompt with independent initial noise. Let $G_i$ and $T_i$
be the student and teacher RGB outputs for the same seed $i$. With frozen
LPIPS-VGG distance $d$, we use
\begin{equation}
L_{\mathrm{div}}=\left[d(G_1,G_2)-\operatorname{sg}(d(T_1,T_2))\right]^2.
\end{equation}
Here $\operatorname{sg}$ stops gradients. The teacher specifies the desired
variation, instead of rewarding unrestricted image separation. Both student
branches receive gradients; teacher, VAE and LPIPS weights remain frozen.
Only the student denoiser and VAE are required during inference. No edge loss
is used in these experiments.
\section{Controlled protocol}
All variants continue the same completed \DMDParentUpdates{}-update DMD parent,
including student/scorer weights, optimizer moments and EMA state. Each receives
\DMDPhaseUpdates{} additional generator and scorer updates, for
\DMDCumulativeUpdates{} cumulative generator updates. Seed is \DMDTrainingSeed{}.
The effective batch is 16, using microbatch two and eight accumulation passes;
the regression stream also contains 16 images per update, grouped by prompt.
All loss settings other than the two experimental factors are matched.

The caption pool consists of \DMDTrainingImages{} fresh COCO-2017 training
images, disjoint from the original 50,000 training images and from validation.
The shared anchor bank contains \DMDPairImages{} generated images for
\DMDPairPrompts{} prompts, two noise seeds per prompt. The frozen SD1.5 teacher
uses 64 DDIM steps and CFG \DMDGuidance{}; pair preprocessing costs
\DMDPairNFE{} denoiser evaluations, performed once and shared across variants.
Regenerating this bank independently for each variant would incur that cost
each time. Generator and scorer learning rates, warmup/cosine G schedule,
weight decay and EMA settings follow the recorded parent configuration.

Each student uses one step and one NFE. FID, Inception Score and CLIP use
\DMDEvalImages{} generated images and the same held-out validation captions,
reference images and evaluation noise. Diversity uses four fixed noise seeds
for each of 128 held-out prompts, averaging six pairwise LPIPS-Alex distances
within each prompt. AlexNet is not used by the training objective. GPU timing
under concurrent workloads is not a controlled speed comparison. LPIPS-VGG
and VAE backpropagation add training cost despite equal generator updates.
\section{Measured results}
\input{generated/dmd_diversity/measured.tex}
\begin{table*}[t]
\centering\small
\input{generated/dmd_diversity/quality.tex}
\caption{Factorial comparison at one step / one NFE. Diversity is LPIPS-Alex
distance across seeds, not an image-quality score. Phase hours exclude the
inherited parent training.}
\end{table*}
\begin{figure*}[t]
\centering
\includegraphics[width=.78\textwidth]{generated/dmd_diversity/samples.pdf}
\caption{The first four consecutive held-out evaluation images. Each row shares
its caption and initial noise across methods; columns identify the four variants.
These samples are selected by fixed order, without quality filtering.}
\end{figure*}
\section{Limitations}
This study has one training seed, reduced prompt coverage in its regression
bank, a rank-32 student and scorer, and a COCO-2017 subset protocol. It does
not reproduce full-model DMD benchmark settings. Matching one perceptual
distance cannot guarantee semantic diversity or prevent every collapsed
solution. A diversity increase may reflect artifacts, so quality and text
alignment must be examined jointly. More training seeds and a fixed-hardware
timing study are required for stronger conclusions. The source code and native
JSON records preserve parent hashes, pair-bank hashes and loss settings.
\bibliographystyle{ieeenat_fullname}
\bibliography{references}
\end{document}
